\documentclass[a4paper]{article}
\usepackage[margin=.8in]{geometry}
\usepackage{amsmath}
\usepackage{amssymb}
\usepackage{enumitem}
\usepackage[cache=true]{minted}
\usepackage{fancyhdr}
\newcommand{\julia}[1]{\mintinline{julia}{#1}}
\setminted{fontsize=\small,autogobble}

\pagestyle{fancy}
\fancyhf{}
\lhead{Name: }
\ifdefined\showsolutions\lhead{Instructor copy}\fi
\chead{Graded quiz 2 -- Introduction to Julia}
\rhead{\today}
\cfoot{\thepage}

\usepackage{xcolor}
\newcommand{\answer}[1]{%
    \ifdefined\showsolutions
        \par\smallskip
        {\small\color{blue!60!black}#1\par}
    \fi
}
% Box in the answer grid, ticked in the solutions when the argument is 1.
\newcommand{\gridbox}[1]{%
    \ifdefined\showsolutions
        \ifnum#1=1 {\color{blue!60!black}$\boxtimes$}\else$\square$\fi
    \else
        $\square$%
    \fi
}

\begin{document}

True or false? (unless otherwise specified.)
Assume any explicitly named package is installed.
An expression that raises an error does not evaluate to \julia{true}.
\ifdefined\showsolutions
{\small\color{blue!60!black}Questions 1--15: 1 point for a correct answer, 0 otherwise.\par}
\fi

\begin{center}
    \textbf{Answer grid}\par\smallskip
    \normalsize
    \renewcommand{\arraystretch}{1.35}
    \begin{tabular}{|c|*{15}{c|}}
        \hline
        & 1 & 2 & 3 & 4 & 5 & 6 & 7 & 8 & 9 & 10 & 11 & 12 & 13 & 14 & 15 \\
        \hline
        True  & \gridbox{1} & \gridbox{1} & \gridbox{1} & \gridbox{1} & \gridbox{0} & \gridbox{1} & \gridbox{1} & \gridbox{1} & \gridbox{0} & \gridbox{1} & \gridbox{0} & \gridbox{1} & \gridbox{0} & \gridbox{1} & \gridbox{0} \\
        False & \gridbox{0} & \gridbox{0} & \gridbox{0} & \gridbox{0} & \gridbox{1} & \gridbox{0} & \gridbox{0} & \gridbox{0} & \gridbox{1} & \gridbox{0} & \gridbox{1} & \gridbox{0} & \gridbox{1} & \gridbox{0} & \gridbox{1} \\
        \hline
    \end{tabular}
\end{center}

\begin{enumerate}[itemsep=8pt]
    \item
        The expression \julia{"ab" * "cd" * "!"} returns the string \julia{"abcd!"}.
        \answer{\textbf{True.}
        \julia{*} concatenates strings.
        The result is \julia{"abcd!"}.}

    \item
        The following code produces a single plot containing both curves,
        rather than replacing the first curve with the second.
        \begin{minted}{julia}
            using Plots
            plot(x -> x^2, xlims=(-1, 1), label="square")
            plot!(x -> x^3, label="cube")
        \end{minted}
        \answer{\textbf{True.}
        \julia{plot!} adds a series to the current plot,
        which already contains the square function.}

    \item
        The expression \julia{eps()/2 + eps()/2 == eps()} evaluates to \julia{true}.
        \answer{\textbf{True.}
        For \julia{Float64}, \julia{eps()} is $2^{-52}$
        and \julia{eps()/2} is $2^{-53}$. Both are exactly representable,
        and adding the two halves gives \julia{eps()} exactly.}

    \item
        The expression \julia{typeof(1 + 2.0) == Float64} evaluates to \julia{true}.
        \answer{\textbf{True.}
        The integer is promoted for this addition;
        the result is \julia{3.0}, of type \julia{Float64}.}

    \item
        If \julia{u = [1, 2]} and \julia{v = [3, 4]}, then \julia{[u v]}
        returns the matrix \julia{[1 2; 3 4]}.
        \answer{\textbf{False.}
        Horizontal concatenation places \julia{u} and
        \julia{v} in the columns, giving \julia{[1 3; 2 4]}.}

    \item
        If \julia{A = [1 2 3; 4 5 6]}, then \julia{A[2, 3]} returns \julia{6}.
        \answer{\textbf{True.}
        The first index selects the row and the second
        selects the column. The entry in row 2, column 3 is 6.}

    \item
        The following code assigns the vector \julia{[2, 5]} to \julia{w}.
        \begin{minted}{julia}
            v = [-3, 2, 0, 5, -1]
            w = v[v .> 0]
        \end{minted}
        \answer{\textbf{True.}
        The mask selects the strictly positive entries,
        giving \julia{[2, 5]}.}

    \item
        The following code assigns the vector \julia{[1, 3, 5]} to \julia{a}.
        \begin{minted}{julia}
            f(x, y) = x - 2y
            a = f.([3, 5, 7], 1)
        \end{minted}
        \answer{\textbf{True.}
        The scalar second argument is used in each call:
        \julia{f(3, 1)}, \julia{f(5, 1)}, and \julia{f(7, 1)}.}

    \item
        The expression \julia{sum(length, ["Julia", "is", "fun"])} returns \julia{3}.
        \answer{\textbf{False.}
        The expression sums the string lengths,
        giving $5+2+3=10$, rather than counting the strings.}

    \item
        The command \julia{A = [i + 2j for i in 1:2, j in 1:3]}
        produces the matrix \julia{A = [3 5 7; 4 6 8]}.
        \answer{\textbf{True.}
        Evaluating \julia{i + 2j} for each row index
        \julia{i} and column index \julia{j} gives \julia{[3 5 7; 4 6 8]}.}

    \item
        After the following definitions, \julia{f(f(1))} returns \julia{2.0}.
        \begin{minted}{julia}
            f(x::Float64) = 1
            f(x::Int) = 2.0
        \end{minted}
        \answer{\textbf{False.}
        The inner call \julia{f(1)} returns \julia{2.0}.
        The outer call therefore uses the \julia{Float64} method and returns \julia{1}.}
\end{enumerate}

\newpage
\begin{enumerate}[resume,itemsep=8pt]
    \item
        The call \julia{odd_sum(5)} returns \julia{9}.
        \begin{minted}{julia}
            function odd_sum(n)
                s = 0
                for k in 1:n
                    if isodd(k)
                        s += k
                    end
                end
                return s
            end
        \end{minted}
        \answer{\textbf{True.}
        The accumulator adds the odd integers $1+3+5=9$.}

    \item
        The expression \julia{LinRange(0, 1, 4) == [0, 0.2, 0.4, 0.6, 0.8, 1]}
        evaluates to \julia{true}.
        \answer{\textbf{False.}
        \julia{LinRange(0, 1, 4)} has four equally spaced
        entries, including both endpoints.
        The vector in the question has six entries instead.}

    \item
        With the following definition, the expression \julia{f(n)} returns
        the $n$-th Fibonacci number for positive integers~$n$
        (assuming no integer overflow).
        Recall that the Fibonacci numbers are defined as
        $f_1=f_2=1$ and $f_k=f_{k-1}+f_{k-2}$ for $k\geq3$.
        \begin{minted}{julia}
            f(n) = n <= 2 ? 1 : f(n-1) + f(n-2)
        \end{minted}
        \answer{\textbf{True.}
        The ternary expression returns 1 for the base cases
        $n=1$ and $n=2$. For $n\geq3$, the recursive calls implement
        the Fibonacci recurrence $f_n=f_{n-1}+f_{n-2}$.}

    \item
        If \julia{v = [2, 4, 6, 8]}, then \julia{v[2:end] == [2, 4, 6]} evaluates to \julia{true}.
        \answer{\textbf{False.}
        \julia{end} denotes the last index, which is 4.
        The slice selects entries 2 through 4, giving \julia{[4, 6, 8]}.}
\end{enumerate}

\noindent\textbf{Bonus (2 points).}
Give the values of \julia{a} and \julia{b} after the following code,
and briefly explain why the results differ.
\begin{minted}{julia}
    v = [[1, 2, 3], [2, 3, 4], [4, 5, 6]]
    a = sum.(v)
    b = sum(v)
\end{minted}
\noindent\textit{Answer and explanation:}
\ifdefined\showsolutions\else\vspace{2cm}\fi
\answer{\julia{a = [6, 9, 15]} and \julia{b = [7, 10, 13]}.
The broadcast \julia{sum.(v)} applies \julia{sum} to each inner vector
independently, returning the vector of the sums $6$, $9$ and $15$.
In contrast, \julia{sum(v)} adds the three vectors together with the
\julia{+} operator, i.e.\ \julia{[1,2,3] + [2,3,4] + [4,5,6]},
yielding the single vector \julia{[7, 10, 13]}.
Award 0.5 points for each correct final value and 0.5 points for each
correct explanation (elementwise sum versus vector addition).}
\end{document}
